% !PRINTED 214
\weberpage{214}
%<<CONT
die sich mit Rücksicht auf $\Theta^2 = \psi^2 = \chi^3 = 1$ aus den drei
Relationen
\[
\psi\Theta = \Theta\psi, \quad \chi\Theta = \Theta\psi\chi, \quad \chi\psi = \Theta\chi
\]
ableiten lassen.

Aus (5) geht noch hervor, dass die Vierergruppe
\weq{6}{1,\ \Theta,\ \psi,\ \Theta\psi}
ein \gesperrt{Normaltheiler} der Tetraëdergruppe ist.

Die Substitutionen der Tetraëdergruppe können in jeder der
beiden Formen
\[
\chi^\lambda \psi^\mu \Theta^\nu, \quad \Theta^\nu \psi^\mu \chi^\lambda
\]
\[
\lambda = 0, 1, 2, \quad \mu = 0, 1, \quad \nu = 0, 1
\]
dargestellt werden.

Um die zu den dreizähligen Polen gehörigen beiden Grundformen
 vierter Ordnung $\Phi_1, \Phi_2$ zu erhalten, kann man entweder
aus (5) die noch fehlenden Pole berechnen, oder man kann so
verfahren:

Da die Grundformen $\Phi_1, \Phi_2$ die Substitutionen $\Theta$ und $\psi$
gestatten müssen und nicht durch $x_1$ und $x_2$ theilbar sein können,
so können sie nur von folgender Form sein:
\[
\begin{aligned}
\Phi_1 &= x_1^4 + m_1 x_1^2 x_2^2 + x_2^4 \\
\Phi_2 &= x_1^4 + m_2 x_1^2 x_2^2 + x_2^4,
\end{aligned}
\]
worin $m_1, m_2$ noch zu bestimmende Constanten sind. Bildet man
aber die \gesperrt{Hesse'sche} Covariante von $\Phi_1$, so erhält man nach
\S.\,40, 4. eine Invariante der Gruppe:
\[
\frac{1}{12}\left[\frac{\partial^2 \Phi_1}{\partial x_1^2}\,\frac{\partial^2 \Phi_1}{\partial x_2^2} - \left(\frac{\partial^2 \Phi_1}{\partial x_1\,\partial x_2}\right)^2\right] = 2m_1(x_1^4+x_2^4) + (12-m_1^2)\,x_1^2 x_2^2.
\]

Diese Form kann, da keine anderen Grundformen vierter
Ordnung existiren, nur entweder mit $\Phi_1$ oder mit $\Phi_2$ übereinstimmen
 (von einem constanten Factor abgesehen). Es muss also
\[
\frac{12-m_1^2}{2m_1} = m_1 \ \text{oder}\ m_2
\]
sein. Die erste Annahme giebt aber $m_1 = \pm 2$. Dann wäre $\Phi_1$
ein Quadrat, was nicht sein kann. Es bleibt also nur:
\[
m_1^2 + 2m_1 m_2 = 12,
\]
und ebenso aus $\Phi_2$:
\[
m_2^2 + 2m_1 m_2 = 12,
\]
also
\[
m_1 = -m_2 = \pm 2\sqrt{-3},
\]
%CONT>>
